\section{AMC 10 10A 2006}
        \begin{enumerate}
        	\item (\textit{AMC 10 10A 2006}) Sandwiches at Joe's Fast Food cost $\textdollar 3$ each and sodas cost $\textdollar 2$ each. How many dollars will it cost to purchase 5 sandwiches and 8 sodas?


 $\mathrm{(A)}\ 31\qquad\mathrm{(B)}\ 32\qquad\mathrm{(C)}\ 33\qquad\mathrm{(D)}\ 34\qquad\mathrm{(E)}\ 35$

 \textbf{Answer: }A



\textbf{Solution}

The $5$ sandwiches cost $5\cdot 3=15$ dollars. The $8$ sodas cost $8\cdot 2=16$ dollars. In total, the purchase costs $15+16=\boxed{\textbf{(A) }31}$ dollars.

 


	\item (\textit{AMC 10 10A 2006}) Define $x\otimes y=x^3-y$. What is $h\otimes (h\otimes h)$?


 $\mathrm{(A)}\ -h\qquad \mathrm{(B)}\ 0\qquad \mathrm{(C)}\ h\qquad \mathrm{(D)}\ 2h\qquad\mathrm{(E)}\ h^3$

 \textbf{Answer: }C



\textbf{Solution 1}

By the definition of $\otimes$, we have $h\otimes h=h^{3}-h$. 

 Then, $h\otimes (h\otimes h)=h\otimes (h^{3}-h)=h^{3}-(h^{3}-h)=\boxed{\textbf{(C) }h}.$



	\item (\textit{AMC 10 10A 2006}) The ratio of Mary's age to Alice's age is $3:5$. Alice is $30$ years old. How many years old is Mary?


 $\mathrm{(A)}\ 15\qquad\mathrm{(B)}\ 18\qquad\mathrm{(C)}\ 20\qquad\mathrm{(D)}\ 24\qquad\mathrm{(E)}\ 50$

 \textbf{Answer: }B



\textbf{Solution 1}

Let $m$ be Mary's age. Then $\frac{m}{30}=\frac{3}{5}$. Solving for $m$, we obtain $m=\boxed{\textbf{(B) }18}.$




\textbf{Solution 2}

We can see this is a combined ratio of $8$, $(5+3)$. We can equalize by doing $30\div5=6$, and $6\cdot3=\boxed{\textbf{(B) }18}$. With the common ratio of $8$ and difference ratio of $6$, we see $6\cdot8=30+18$. Therefore, we can see our answer is correct.

 


	\item (\textit{AMC 10 10A 2006}) A digital watch displays hours and minutes with AM and PM.  What is the largest possible sum of the digits in the display?


 $\mathrm{(A)}\ 17\qquad\mathrm{(B)}\ 19\qquad\mathrm{(C)}\ 21\qquad\mathrm{(D)}\ 22\qquad\mathrm{(E)}\ 23$

 \textbf{Answer: }E



\textbf{Solution 1}

From the \href{https://artofproblemsolving.com/wiki/index.php?title=Greedy\_algorithm}{greedy algorithm}, we have $9$ in the hours section and $59$ in the minutes section. $9+5+9=\boxed{\textbf{(E) }23}$




\textbf{Solution 2 (matrix)}

With a matrix, we can see$\begin{bmatrix} 1+2&9&6&3\\ 1+1&8&5&2\\ 1+0&7&4&1 \end{bmatrix}$The largest single digit sum we can get is $9$.For the minutes digits, we can combine the largest $2$ digits, which are $9,5 \Rightarrow 9+5=14$, and finally $14+9=\boxed{\textbf{(E) }23}$




\textbf{Solution 3}

We first note that since the watch displays time in AM and PM, the value for the hours section varies from $00-12$. Therefore, the maximum value of the digits for the hours is when the watch displays $09$, which gives us $0+9=9$.

 Next, we look at the value of the minutes section, which varies from $00-59$. Let this value be a number $ab$. We quickly find that the maximum value for $a$ and $b$ is respectively $5$ and $9$.

 Adding these up, we get $9+5+9=\boxed{\textbf{(E) }23}$.

 


	\item (\textit{AMC 10 10A 2006}) Doug and Dave shared a pizza with 8 equally-sized slices.  Doug wanted a plain pizza, but Dave wanted anchovies on half of the pizza.  The cost of a plain pizza was 8 dollars, and there was an additional cost of 2 dollars for putting anchovies on one half.  Dave ate all of the slices of anchovy pizza and one plain slice.  Doug ate the remainder.  Each then paid for what he had eaten.  How many more dollars did Dave pay than Doug?


 $\mathrm{(A)}\ 1\qquad\mathrm{(B)}\ 2\qquad\mathrm{(C)}\ 3\qquad\mathrm{(D)}\ 4\qquad\mathrm{(E)}\ 5$

 \textbf{Answer: }D



\textbf{Solution}

Dave and Doug paid $8+2=10$ dollars in total.  Doug paid for three slices of plain pizza, which cost $\frac{3}{8}\cdot 8=3$.  Dave paid $10-3=7$ dollars.  Dave paid $7-3=\boxed{\textbf{(D) }4}$ more dollars than Doug.

 


	\item (\textit{AMC 10 10A 2006}) What non-zero real value for $x$ satisfies $(7x)^{14}=(14x)^7$?


 $\mathrm{(A)}\ \frac17\qquad\mathrm{(B)}\ \frac27\qquad\mathrm{(C)}\ 1\qquad\mathrm{(D)}\ 7\qquad\mathrm{(E)}\ 14$

 \textbf{Answer: }B



\textbf{Solution}

Taking the seventh root of both sides, we get $(7x)^2=14x$.

 Simplifying the LHS gives $49x^2=14x$. This means that $7x = 2$, or $x = 0$.

 However, $x$ must be a non-zero value.

 Thus, $x=\boxed{\textbf{(B) }\frac{2}{7}}$.

 


	\item (\textit{AMC 10 10A 2006}) 
                    \begin{figure}[H]
                        \centering
                        \includegraphics[width=0.5\textwidth]{Problem/AMC_10_10A_2006_7_Problem_Figure_1.png}
                    \end{figure}
                    


 The $8 \times 18$ rectangle $ABCD$ is cut into two congruent hexagons, as shown, in such a way that the two hexagons can be repositioned without overlap to form a square.  What is $y$?


 $\mathrm{(A)}\ 6\qquad\mathrm{(B)}\ 7\qquad\mathrm{(C)}\ 8\qquad\mathrm{(D)}\ 9\qquad\mathrm{(E)}\ 10$

 \textbf{Answer: }A



\textbf{Solution 1}

Since the two \href{https://artofproblemsolving.com/wiki/index.php?title=Hexagon}{hexagons} are going to be repositioned to form a \href{https://artofproblemsolving.com/wiki/index.php?title=Square\_(geometry)}{ square} without overlap, the \href{https://artofproblemsolving.com/wiki/index.php?title=Area}{area} will remain the same. The rectangle's area is $18\cdot8=144$. This means the square will have four sides of length 12. The only way to do this is shown below.


 
                    \begin{figure}[H]
                        \centering
                        \includegraphics[width=0.5\textwidth]{Solution/AMC_10_10A_2006_7_Solution_Figure_1.png}
                    \end{figure}
                    

 As you can see from the diagram, the \href{https://artofproblemsolving.com/wiki/index.php?title=Line\_segment}{line segment} denoted as $y$ is half the length of the side of the square, which leads to $y = \frac{12}{2} = \boxed{\textbf{(A) }6}$.

 



\textbf{Solution 2  (Shortcut)}

As solution 1 says, the two hexagons are going to be repositioned to form a square without overlap. Thus we create this square out of the original rectangle.

 

                    \begin{figure}[H]
                        \centering
                        \includegraphics[width=0.5\textwidth]{Solution/AMC_10_10A_2006_7_Solution_Figure_2.png}
                    \end{figure}
                    

 As you can see from the diagram, the length $y$ fits into the previously blank side, so we know that it is equal to $y$.

 
                    \begin{figure}[H]
                        \centering
                        \includegraphics[width=0.5\textwidth]{Solution/AMC_10_10A_2006_7_Solution_Figure_3.png}
                    \end{figure}
                    

 
From there we can say $3y = 18$ so $y = \frac{18}{3} = \boxed{\textbf{(A) }6}$.

 



\textbf{Solution 3 (Cheap)}

Because the two hexagons are congruent, we know that the perpendicular line to $A$ is half of $BC$, or $4$. Next, we plug the answer choices in to see which one works. Trying $A$, we get the area of one hexagon is $72$ , as desired, so the answer is $\boxed{\textbf{(A) }6}$ .

 


	\item (\textit{AMC 10 10A 2006}) A parabola with equation $y=x^2+bx+c$ passes through the points $(2,3)$ and $(4,3)$.  What is $c$?


 $\mathrm{(A)}\ 2\qquad\mathrm{(B)}\ 5\qquad\mathrm{(C)}\ 7\qquad\mathrm{(D)}\ 10\qquad\mathrm{(E)}\ 11$

 \textbf{Answer: }E



\textbf{Solution 1}

Substitute the points $(2,3)$ and $(4,3)$ into the given equation for $(x,y)$.

 Then we get a system of two equations:

 $3=4+2b+c$
$3=16+4b+c$

 Subtracting the first equation from the second we have:

 $0=12+2b$
$b=-6$

 Then using $b=-6$ in the first equation:

 $0=1+-12+c$
$c=\boxed{\textbf{(E) }11}$.

 



\textbf{Solution 1.1}

Alternatively, notice that since the equation is that of a conic parabola, the vertex is likely $(3,2)$. Thus, the form of the equation of the parabola is $y - 2 = (x - 3)^2$. Expanding this out, we find that $c=\boxed{\textbf{(E) }11}$.

 



\textbf{Solution 2}

The points given have the same $y$-value, so the vertex lies on the line $x=\frac{2+4}{2}=3$.

 The $x$-coordinate of the vertex is also equal to $\frac{-b}{2a}$, so set this equal to $3$ and solve for $b$, given that $a=1$:

 $x=\frac{-b}{2a}$
$3=\frac{-b}{2}$
$6=-b$
$b=-6$

 Now the equation is of the form $y=x^2-6x+c$.  Now plug in the point $(2,3)$ and solve for $c$:

 $y=x^2-6x+c$
$3=2^2-6(2)+c$
$3=4-12+c$
$3=-8+c$
$c=\boxed{\textbf{(E) }11}$.

 



\textbf{Solution 3}

Substituting y into the two equations, we get:

 $3=x^2+bx+c$

 Which can be written as:

 $x^2+bx+c-3=0$
$4$ and $2$ are the solutions to the quadratic. Thus:

 $c-3=4\times2$
$c-3=8$
$c=\boxed{\textbf{(E) }11}$.

 


	\item (\textit{AMC 10 10A 2006}) How many sets of two or more consecutive positive integers have a sum of 15?


 $\mathrm{(A)}\ 1\qquad\mathrm{(B)}\ 2\qquad\mathrm{(C)}\ 3\qquad\mathrm{(D)}\ 4\qquad\mathrm{(E)}\ 5$

 \textbf{Answer: }C



\textbf{Solution 1}

Notice that if the consecutive positive integers have a sum of $15$, then their average (which could be a fraction) must be a divisor of $15$. If the number of integers in the list is odd, then the average must be either $1, 3,$ or $5$, and $1$ is clearly not possible. The other two possibilities both work:

 
            \begin{itemize}
        \item $1 + 2 + 3 + 4 + 5 = 15$
        \item $4 + 5 + 6 = 15$
            \end{itemize}
            If the number of integers in the list is even, then the average will have a $\frac{1}{2}$. The only possibility is $\frac{15}{2}$, from which we get:

 
            \begin{itemize}
        \item $15 = 7 + 8$
            \end{itemize}
            Thus, the correct answer is $\boxed{\textbf{(C) }3}.$

 
Question: (RealityWrites) Is it possible that the answer is $4$, because $0+1+2+3+4+5$ should technically count, right?

 Answer: (IMGROOT2) It isn't possible because the question asks for positive integers, and this means that negative integers or zero aren't allowed.

 Note to readers: make sure to always read the problem VERY carefully before attempting; it could mean the difference of making the cutoff.

 



\textbf{Solution 2}

Any set will form a arithmetic progression with the first term say $a$. Since the numbers are consecutive the common difference $d = 1$.

 The sum of the AP has to be 15. So,

 
 \[S_n = \frac{n}{2} \cdot (2a + (n-1)d)\]
 \[S_n = \frac{n}{2} \cdot (2a + (n-1)1)\]
 \[15 = \frac{n}{2} \cdot (2a + n - 1)\]
 \[2an + n^2 - n = 30\]
 \[n^2 + n(2a - 1) - 30 = 0\]

 Now we need to find the number of possible sets of values of a, n which satisfy this equation. Now $a$ cannot be 15 as we need 2 terms. So a can only be less the 15.

 Trying all the values of a from 1 to 14 we observe that $a = 4$, $a = 7$ and $a = 1$ provide the only real solutions to the above equation.The three possibilites of a and n are.

 
 \[(a,n) = (4,3),(7, 2),(1, 6)\]

 The above values are obtained by solving the following equations obtained by substituting the above mentioned values of a into $n^2 + n(2a - 1) - 30 = 0$,

 
 \[n^2 + 7n - 30 = 0\]
 \[n^2 + 13n - 30 = 0\]
 \[n^2 - n - 30 = 0\]

 Since there are 3 possibilities the answer is $\boxed{\textbf{(C) }3}.$




\textbf{Solution 3}

We want some consecutive numbers that add up. We can use casework to solve this. We know that any set of 6 consecutive, positive integers cannot be greater than 5. Therefore the max amount of elements is 5.

 \textbf{Case 1: There are 5 numbers} -We know that 15/5 = 3, so if we put 3 as our middle number and expand evenly, we get {1,2,3,4,5}

 \textbf{Case 2: There are 4 numbers} -By a bit of bashing we find that there are no solutions

 \textbf{Case 3: There are 3 numbers} -We know that 15/3 = 5, so if we put 5 as our middle number and expand evenly, we get {4,5,6}

 \textbf{Case 4: There are 2 numbers} -We'll say x is the lower number. Thus x + (x + 1) = 15. Solving for x, x = 7. Thus getting {7,8}

 Adding up all the cases we get 1 + 0 + 1 + 1 = $\boxed{\textbf{(C) }3}.$



	\item (\textit{AMC 10 10A 2006}) For how many real values of $x$ is $\sqrt{120-\sqrt{x}}$ an integer?


 $\mathrm{(A)}\ 3\qquad\mathrm{(B)}\ 6\qquad\mathrm{(C)}\ 9\qquad\mathrm{(D)}\ 10\qquad\mathrm{(E)}\ 11\qquad$

 \textbf{Answer: }E



\textbf{Solution}

For $\sqrt{120-\sqrt{x}}$ to be an integer, $120-\sqrt{x}$ must be a perfect square. 

 Since $\sqrt{x}$ can't be negative, $120-\sqrt{x} \leq 120$. 

 The perfect squares that are less than or equal to $120$ are $\{0,1,4,9,16,25,36,49,64,81,100\}$, so there are $11$ values for $120-\sqrt{x}$. 

 Since every value of $120-\sqrt{x}$ gives one and only one possible value for $x$, the number of values of $x$ is $\boxed{\textbf{(E) }11}$.

 


	\item (\textit{AMC 10 10A 2006}) Which of the following describes the graph of the equation $(x+y)^2=x^2+y^2$?


 $\mathrm{(A)}\ \text{the empty set}\qquad\mathrm{(B)}\ \text{one point}\qquad\mathrm{(C)}\ \text{two lines}\qquad\mathrm{(D)}\ \text{a circle}\qquad\mathrm{(E)}\ \text{the entire plane}$

 \textbf{Answer: }C



\textbf{Solution}

Expanding the left side, we have

 $x^2+2xy+y^2=x^2+y^2\Longrightarrow 2xy=0\Longrightarrow xy=0\Longrightarrow x = 0 \textrm{ or } y = 0$

 Thus there are two \href{https://artofproblemsolving.com/wiki/index.php?title=Line}{lines} described in this graph, the horizontal line $y = 0$ and the vertical line $x=0$. Thus, our answer is $\boxed{\textbf{(C) }\textrm{two\,lines}}$.

 


	\item (\textit{AMC 10 10A 2006}) 
                    \begin{figure}[H]
                        \centering
                        \includegraphics[width=0.5\textwidth]{Problem/AMC_10_10A_2006_12_Problem_Figure_1.png}
                    \end{figure}
                    


 Rolly wishes to secure his dog with an 8-foot rope to a square shed that is 16 feet on each side.  His preliminary drawings are shown.


 Which of these arrangements give the dog the greater area to roam, and by how many square feet?


 $\mathrm{(A)}\ \text{I, by}\ 8\pi\qquad\mathrm{(B)}\ \text{I, by}\ 6\pi\qquad\mathrm{(C)}\ \text{II, by}\ 4\pi\qquad\mathrm{(D)} \text{II, by}\ 8\pi\qquad\mathrm{(E)}\ \text{II, by}\ 10\pi$

 \textbf{Answer: }C



\textbf{Solution}


                    \begin{figure}[H]
                        \centering
                        \includegraphics[width=0.5\textwidth]{Solution/AMC_10_10A_2006_12_Solution_Figure_1.png}
                    \end{figure}
                    Let us first examine the area of both possible arrangements. The rope outlines a circular boundary that the dog may dwell in. Arrangement $I$ allows the dog $\frac12\cdot(\pi\cdot8^2) = 32\pi$ square feet of area. Arrangement $II$ allows $32\pi$ square feet plus a little more on the top part of the fence. So we already know that Arrangement $II$ allows more freedom - only thing left is to find out how much. The extra area can be represented by a quarter of a circle with radius 4. So the extra area is $\frac14\cdot(\pi\cdot4^2) = 4\pi$. Thus the answer is $\boxed{\textbf{(C) } II,\,\textrm{ by }\,4\pi}$.

 


	\item (\textit{AMC 10 10A 2006}) A player pays $\textdollar 5$ to play a game. A die is rolled. If the number on the die is odd, the game is lost. If the number on the die is even, the die is rolled again. In this case the player wins if the second number matches the first and loses otherwise. How much should the player win if the game is fair? (In a fair game the probability of winning times the amount won is what the player should pay.) 


 $\mathrm{(A)}\ 12\qquad\mathrm{(B)}\  30\qquad\mathrm{(C)}\  50\qquad\mathrm{(D)}\  60\qquad\mathrm{(E)}\  100$

 \textbf{Answer: }D



\textbf{Solution}

The probability of rolling an even number on the first turn is $\frac{1}{2}$ and the probability of rolling the same number on the next turn is $\frac{1}{6}$. The probability of winning is $\frac{1}{2}\cdot \frac{1}{6} =\frac{1}{12}$. If the game is to be fair, the amount paid, $5$ dollars, must be $\frac{1}{12}$ the amount of the prize money, so the answer is$\boxed{\textbf{(D) } \$60}.$



	\item (\textit{AMC 10 10A 2006}) 
                    \begin{figure}[H]
                        \centering
                        \includegraphics[width=0.5\textwidth]{Problem/AMC_10_10A_2006_14_Problem_Figure_1.png}
                    \end{figure}
                    


 A number of linked rings, each 1 cm thick, are hanging on a peg. The top ring has an outside diameter of 20 cm. The outside diameter of each of the other rings is 1 cm less than that of the ring above it. The bottom ring has an outside diameter of 3 cm. What is the distance, in cm, from the top of the top ring to the bottom of the bottom ring? 


 $\mathrm{(A)}\ 171\qquad\mathrm{(B)}\ 173\qquad\mathrm{(C)}\ 182\qquad\mathrm{(D)}\ 188\qquad\mathrm{(E)}\ 210$

 \textbf{Answer: }B



\textbf{Solution 1}

The inside diameters of the rings are the positive integers from $1$ to $18$.  The total distance needed is the sum of these values plus $2$ for the top of the first ring and the bottom of the last ring. Using the formula for the sum of an \href{https://artofproblemsolving.com/wiki/index.php?title=Arithmetic\_series}{arithmetic series}, the answer is $\frac{18 \cdot 19}{2} + 2 = \boxed{\textbf{(B) }173}$.

 



\textbf{Solution 2}

Alternatively, the sum of the consecutive integers from 3 to 20 is $\frac{1}{2}(18)(3+20) = 207$.  However, the 17 \href{https://artofproblemsolving.com/wiki/index.php?title=Intersection}{intersections} between the rings must be subtracted, and we also get $207 - 2(17) = \boxed{\textbf{(B) }173}$.

 



\textbf{Solution 3 (Brute force)}

We can just add all the numbers up. But, we have to account for the intersections, so the final expression is $3+\sum_{k=2}^{18} k$, which is $\boxed{\textbf{(B) }173}$. ~\href{https://artofproblemsolving.com/wiki/index.php?title=Shunyipanda&action=edit&redlink=1}{shunyipanda}

 


	\item (\textit{AMC 10 10A 2006}) Odell and Kershaw run for 30 minutes on a circular track. Odell runs clockwise at 250 m/min and uses the inner lane with a radius of 50 meters. Kershaw runs counterclockwise at 300 m/min and uses the outer lane with a radius of 60 meters, starting on the same radial line as Odell. How many times after the start do they pass each other? 


 $\mathrm{(A)}\ 29\qquad\mathrm{(B)}\ 42\qquad\mathrm{(C)}\ 45\qquad\mathrm{(D)}\ 47\qquad\mathrm{(E)}\ 50$

 \textbf{Answer: }D



\textbf{Solution 1}


                    \begin{figure}[H]
                        \centering
                        \includegraphics[width=0.5\textwidth]{Solution/AMC_10_10A_2006_15_Solution_Figure_1.png}
                    \end{figure}
                    

 Since $d = rt$, we note that Odell runs one lap in $\frac{2 \cdot 50\pi}{250} = \frac{2\pi}{5}$ minutes, while Kershaw also runs one lap in $\frac{2 \cdot 60\pi}{300} = \frac{2\pi}{5}$ minutes. They take the same amount of time to run a lap, and since they are running in opposite directions they will meet exactly twice per lap (once at the starting point, the other at the half-way point). Thus, there are $\frac{30}{\frac{2\pi}{5}} \approx 23.8$ laps run by both, or $\lfloor 2\cdot 23.8\rfloor = 23 \cdot 2 + 1 =\boxed{\textbf{(D) } 47}$ meeting points.

 



\textbf{Solution 2}

We first find the amount of minutes, $k$, until Odell and Kershaw's next meeting. Let $a$ be the angle in \href{https://artofproblemsolving.com/wiki/index.php?title=Radians}{radians} between their starting point and the point where they first meet, measured counterclockwise. 

 Since Kershaw has traveled $300k$ meters at this point and the circumference of his track is $120\pi$, $a=\frac{300k}{120\pi}\cdot 2\pi$. Similarly, $2\pi-a=\frac{250k}{100\pi}\cdot{2\pi}$ since Odell has traveled $250k$ meters in the opposite direction and the circumference of his track is $100\pi$. 

 Solving for $a$ in the second equation, we get $a=2\pi-\frac{250k}{100\pi}\cdot 2\pi$. Then, from the first equation, we have $\frac{300k}{120\pi}\cdot 2\pi=2\pi-\frac{250k}{100\pi}\cdot 2\pi$. Solving for $k$, we get $k=\frac{\pi}{5}$. After $k$ minutes, they are back at the same position, except rotated, so they will meet again in $k$ minutes. So the total amount of meetings is $\lfloor\frac{30}{k}\rfloor=\lfloor\frac{150}{\pi}\rfloor=\boxed{\textbf{(D) }47}$.

 



\textbf{Solution 3 by Alcumus (ikr)}

Since Odell's rate is $5/6$ that of Kershaw, but Kershaw's lap distance is $6/5$ that of Odell, they each run a lap in the same time. Hence they pass twice each time they circle the track. Odell runs 
 \[(30 \ \text{min})\left(250\frac{\text{m}}{\text{min}}\right)\left(\frac{1}{100\pi}\frac{\text{laps}}{\text{m}}\right)= \frac{75}{\pi}\,\text{laps}\approx 23.87\ \text{laps},\] as does Kershaw. Because $23.5 < 23.87 < 24$, they pass each other $2(23.5)=\boxed{47}$ times.

 


	\item (\textit{AMC 10 10A 2006}) 
                    \begin{figure}[H]
                        \centering
                        \includegraphics[width=0.5\textwidth]{Problem/AMC_10_10A_2006_16_Problem_Figure_1.png}
                    \end{figure}
                    


 A circle of radius 1 is tangent to a circle of radius 2. The sides of $\triangle ABC$ are tangent to the circles as shown, and the sides $\overline{AB}$  and $\overline{AC}$  are congruent. What is the area of $\triangle ABC$?


 $\mathrm{(A)}\ \frac{35}{2}\qquad\mathrm{(B)}\ 15\sqrt{2}\qquad\mathrm{(C)}\ \frac{64}{3}\qquad\mathrm{(D)}\ 16\sqrt{2}\qquad\mathrm{(E)}\ 24$

 \textbf{Answer: }D



\textbf{Solution 1}

Let the centers of the smaller and larger circles be $O_1$ and $O_2$ , respectively.Let their tangent points to $\triangle ABC$ be $D$ and $E$, respectively.We can then draw the following diagram:

 
                    \begin{figure}[H]
                        \centering
                        \includegraphics[width=0.5\textwidth]{Solution/AMC_10_10A_2006_16_Solution_Figure_1.png}
                    \end{figure}
                    

 We see that $\triangle ADO_1 \sim \triangle AEO_2 \sim \triangle AFC$. Using the first pair of \href{https://artofproblemsolving.com/wiki/index.php?title=Similar\_triangles}{similar triangles}, we write the proportion:

 $\frac{AO_1}{AO_2} = \frac{DO_1}{EO_2} \Longrightarrow \frac{AO_1}{AO_1 + 3} = \frac{1}{2} \Longrightarrow AO_1 = 3$By the \href{https://artofproblemsolving.com/wiki/index.php?title=Pythagorean\_Theorem}{Pythagorean Theorem}, we have $AD = \sqrt{3^2-1^2} = \sqrt{8}$.

 Now using $\triangle ADO_1 \sim \triangle AFC$,

 $\frac{AD}{AF} = \frac{DO_1}{FC} \Longrightarrow \frac{2\sqrt{2}}{8} = \frac{1}{FC} \Longrightarrow FC = 2\sqrt{2}$Hence, the area of the triangle is 
 \[\frac{1}{2}\cdot AF \cdot BC = \frac{1}{2}\cdot AF \cdot (2\cdot CF) = AF \cdot CF = 8\left(2\sqrt{2}\right) = \boxed{\textbf{(D) } 16\sqrt{2}}\]




\textbf{Solution 2}


                    \begin{figure}[H]
                        \centering
                        \includegraphics[width=0.5\textwidth]{Solution/AMC_10_10A_2006_16_Solution_Figure_1.png}
                    \end{figure}
                    

 Since $\triangle{A O_1 D} \sim \triangle{A O_2 E},$ we have that $\frac{A O_1}{A O_2} = \frac{O_1 D}{O_2 E} = \frac{1}{2}$.

 Since we know that $O_1 O_2 = 1 + 2 = 3,$ the total length of $A O_2 = 2 \cdot 3 = 6.$

 We also know that $O_2 F = 2$, so $A F = A O_2 + O_2 F = 6 + 2 = 8.$

 Also, since $\triangle{ABF} \sim \triangle{A E O_2},$ we have that $\frac{AC}{A O_2} = \frac{FC}{O_2 E}.$

 Since we know that $A O_2 = 6$ and $O_2 E = 2,$ we have that $\frac{AC}{6} = \frac{FC}{2}.$

 This equation simplified gets us $AC = 3 \cdot FC.$

 Let $FC = a$

 By the \href{https://artofproblemsolving.com/wiki/index.php?title=Pythagorean\_Theorem}{Pythagorean Theorem} on $\triangle{AFC},$ we have that $AF^2 + FC^2 = AC^2.$

 We know that $AF = 8$, $FC = a$ and $AC = 3a$ so we have $8^2 + a^2 = (3a)^2.$

 Simplifying, we have that $64 = 8a^2 \Rightarrow a^2 = 8 \Rightarrow a = \sqrt{8} = 2 \sqrt{2}.$

 Recall that $FC=a$.

 Therefore, $BC = 2 \cdot FC = 2 \cdot 2 \sqrt{2} = 4 \sqrt{2}.$

 Since the height is $AF = 8,$ we have the area equal to $\frac{4 \sqrt{2} \cdot 8}{2}=16\sqrt{2}.$

 Thus our answer is $\boxed{\textbf{(D) }16 \sqrt{2}}$.

 



\textbf{Solution 3}

Using the diagram from Solutions 1 and 2, and knowing that the height has length eight, we can use the Pythagorean Theorem on triangle $AFC$. We get that 
 \[AC = \sqrt{FC^2 + 64}.\] We can set up the equation 
 \[2FC + 2 \sqrt{FC^2 + 64} = 8FC,\] both sides equal the full perimeter of the triangle. Subtracting $2FC$ from both sides gives 
 \[2\sqrt{FC^2+64}=6FC.\] Squaring and dividing both sides by four gives $FC^2 + 64 = 9FC^2$, meaning that $64 = 8FC^2.$ This implies $FC = 2 \sqrt{2}$. The whole length $BC$ is $4\sqrt{2}$. Thus, the area of the triangle is $16\sqrt{2}\Rightarrow\boxed{D}.$ ~booking~edits by HM2018

 


	\item (\textit{AMC 10 10A 2006}) In rectangle $ADEH$, points $B$ and $C$ trisect $\overline{AD}$, and points $G$ and $F$ trisect $\overline{HE}$. In addition, $AH=AC=2$, and $AD=3$. What is the area of quadrilateral $WXYZ$ shown in the figure? 


 
                    \begin{figure}[H]
                        \centering
                        \includegraphics[width=0.5\textwidth]{Problem/AMC_10_10A_2006_17_Problem_Figure_1.png}
                    \end{figure}
                    
$\textbf{(A) } \frac{1}{2}\qquad\textbf{(B) } \frac{\sqrt{2}}{2}\qquad\textbf{(C) } \frac{\sqrt{3}}{2}\qquad\textbf{(D) } \sqrt{2} \qquad\textbf{(E) } \frac{2\sqrt{3}}{3}\qquad$

 \textbf{Answer: }A



\textbf{Solution 1}

By \href{https://artofproblemsolving.com/wiki/index.php?title=Symmetry}{symmetry}, $WXYZ$ is a square.

 
                    \begin{figure}[H]
                        \centering
                        \includegraphics[width=0.5\textwidth]{Solution/AMC_10_10A_2006_17_Solution_Figure_1.png}
                    \end{figure}
                    

 Draw 

 


	\item (\textit{AMC 10 10A 2006}) A license plate in a certain state consists of 4 digits, not necessarily distinct, and 2 letters, also not necessarily distinct. These six characters may appear in any order, except that the two letters must appear next to each other. How many distinct license plates are possible? 


 $\mathrm{(A)}\ 10^4\times26^2\qquad\mathrm{(B)}\ 10^3\times26^3\qquad\mathrm{(C)}\ 5\times10^4\times26^2\qquad\mathrm{(D)}\ 10^2\times26^4\qquad\mathrm{(E)}\ 5\times10^3\times26^3$

 \textbf{Answer: }C



\textbf{Solution}

There are $10\cdot10\cdot10\cdot10 = 10^4$ ways to choose 4 digits.

 There are $26 \cdot 26 = 26^2$ ways to choose the 2 letters. 

 For the letters to be next to each other, they can be the 1st and 2nd, 2nd and 3rd, 3rd and 4th, 4th and 5th, or the 5th and 6th characters. So, there are $6 - 1 = 5$ choices for the position of the letters.

 Therefore, the number of distinct license plates is  $5\times 10^4\times 26^2 \Longrightarrow \boxed{\mathrm{C}}$.

 



\textbf{Solution 2 (Answer choices)}

There are $10^4$ ways to choose the digits and $26^2$ ways to choose the 2 letters. However, since the letters are next to each other, the result must be $10^4\cdot 26^2 \cdot \text{something}$. The only answer choice that matches this is $\boxed{\mathrm{C}}$.

 


	\item (\textit{AMC 10 10A 2006}) How many non-similar triangles have angles whose degree measures are distinct positive integers in arithmetic progression? 


 $\mathrm{(A)}\ 0\qquad\mathrm{(B)}\ 1\qquad\mathrm{(C)}\ 59\qquad\mathrm{(D)}\ 89\qquad\mathrm{(E)}\ 178$

 \textbf{Answer: }C



\textbf{Solution}

The sum of the angles of a triangle is $180$ degrees. For an arithmetic progression with an odd number of terms, the middle term is equal to the average of the sum of all of the terms, making it $\frac{180}{3} = 60$ degrees. The minimum possible value for the smallest angle is $1$ and the highest possible is $59$ (since the numbers are distinct), so there are $\boxed{\textbf{(C) }59}$ possibilities.

 



\textbf{Solution 2 (Stars and Bars)}

Let the first angle be $x$, and the common difference be $d$. The arithmetic progression can now be expressed as $x + (x + d) + (x + 2d) = 180$. Simplifiying, $x + d = 60$. Now, using stars and bars, we have $\binom{61}{1} = 61$. However, we must subtract the two cases in which either $x$ or $d$ equal $0$, so we have $61 - 2$ = $\boxed{\textbf{(C) }59}$.

 



\textbf{Solution 3 (Quick Summation)}

Consider that we have $(a+n)+(a+n+1)+(a+n+2)=180 \Longleftrightarrow 3a+3(n+1)=180 \Longleftrightarrow a=59-n$, where $n \geq 0$ and $n$ is an integer. Since $a \neq 0$, $n=0,1,2,3,\cdots, 58$ which is $\boxed{\textbf{(C) }59}$ solutions.

 


	\item (\textit{AMC 10 10A 2006}) Six distinct positive integers are randomly chosen between 1 and 2006, inclusive. What is the probability that some pair of these integers has a difference that is a multiple of 5? 


 $\mathrm{(A)}\ \frac{1}{2}\qquad\mathrm{(B)}\ \frac{3}{5}\qquad\mathrm{(C)}\ \frac{2}{3}\qquad\mathrm{(D)}\ \frac{4}{5}\qquad\mathrm{(E)}\ 1$

 \textbf{Answer: }E



\textbf{Solution}

For two numbers to have a difference that is a multiple of $5$, the numbers must be congruent $\bmod{5}$ (their remainders after division by $5$ must be the same).

 $0, 1, 2, 3, 4$ are the possible values of numbers in $\bmod{5}$. Since there are only $5$ possible values in $\bmod{5}$ and we are picking $6$ numbers, by the \href{https://artofproblemsolving.com/wiki/index.php?title=Pigeonhole\_Principle}{Pigeonhole Principle}, two of the numbers must be congruent $\bmod{5}$. 

 Therefore the probability that some pair of the $6$ integers has a difference that is a multiple of $5$ is $\boxed{\textbf{(E) }1}$.

 


	\item (\textit{AMC 10 10A 2006}) How many four-\href{https://artofproblemsolving.com/wiki/index.php?title=Digit}{digit} \href{https://artofproblemsolving.com/wiki/index.php?title=Positive\_integer}{positive integers} have at least one digit that is a $2$ or a $3$?


 $\textbf{(A) } 2439\qquad\textbf{(B) } 4096\qquad\textbf{(C) } 4903\qquad\textbf{(D) } 4904\qquad\textbf{(E) } 5416$

 \textbf{Answer: }E



\textbf{Solution 1 (Complementary Counting)}

Since we are asked for the number of positive $4$-digit integers with at least $2$ or $3$ in it, we can find this by finding the total number of $4$-digit integers and subtracting off those which do not have any $2$s or $3$s as digits. 

 The total number of $4$-digit integers is $9 \cdot 10 \cdot 10 \cdot 10 = 9000$, since we have $10$ choices for each digit except the first (which can't be $0$).

 Similarly, the total number of $4$-digit integers without any $2$ or $3$ is $7 \cdot 8 \cdot 8 \cdot 8 ={3584}$.

 Therefore, the total number of positive $4$-digit integers that have at least one $2$ or $3$ is $9000-3584=\boxed{\textbf{(E) }5416}.$




\textbf{Solution 2 (Casework)}

We proceed to every case.

 Case $1$: There is ONLY one $2$ or $3$. If the $2$ or $3$ is occupying the first digit, we have $512$ arrangements. If the $2$ or $3$ is not occupying the first digit, there are $7 \cdot 8^2$ = $448$ arrangements. Therefore, we have $2(448 \cdot 3 + 512) = 3712$ arrangements.

 Case $2$ : There are two $2$s OR two $3$s. If the $2$ or $3$ is occupying the first digit, we have $64$ arrangements. If the $2$ or $3$ is not occupying the first digit, there are $56$ arrangements. There are $3$ ways for the $2$ or the $3$ to be occupying the first digit and $3$ ways for the first digit to be unoccupied. There are $2(3 \cdot (56+64))$ = $720$ arrangements.

 Case $3$ : There is ONLY one $2$ and one $3$. If the $2$ or the $3$ is occupying the first digit, we have $6$ types of arrangements of where the $2$ or $3$ is. We also have $64$ different arrangements for the non-$2$ or $3$ digits. We have $6 \cdot 64$ = $384$ arrangements. If the $2$ or the $3$ isn't occupying the first digit, we have $6$ types of arrangements of where the $2$ or $3$ is. We also have $56$ different arrangements for the non-$2$ or $3$ digits. We have $6 \cdot 56$ = $336$ arrangements for this case. We have $336 + 384$ = $720$ total arrangements for this case.

 Notice that we already counted $3712 + 720 + 720 = 5152$ cases and we still have a lot of cases left over to count. This is already larger than the second largest answer choice, and therefore, our answer is $\boxed{\textbf{(E) }5416}$.

 


	\item (\textit{AMC 10 10A 2006}) Two farmers agree that pigs are worth $300 and that goats are worth $210. When one farmer owes the other money, he pays the debt in pigs or goats, with "change" received in the form of goats or pigs as necessary. (For example, a $390 debt could be paid with two pigs, with one goat received in change.) What is the amount of the smallest positive debt that can be resolved in this way? 


 $\mathrm{(A)}\ 5\qquad\mathrm{(B)}\ 10\qquad\mathrm{(C)}\ 30\qquad\mathrm{(D)}\ 90\qquad\mathrm{(E)}\ 210$

 \textbf{Answer: }C



\textbf{Solution 1 (Diophantine Equation)}

The problem can be restated as an equation of the form $300p + 210g = x$, where $p$ is the number of pigs, $g$ is the number of goats, and $x$ is the positive debt. The problem asks us to find the lowest \textit{x} possible. \textit{$p$} and \textit{$g$} must be integers, which makes the equation a \href{https://artofproblemsolving.com/wiki/index.php?title=Diophantine\_equation}{Diophantine equation}.  \href{https://artofproblemsolving.com/wiki/index.php?title=Bezout\%27s\_Lemma}{Bezout's Lemma} tells us that the smallest $c$ for the Diophantine equation $am + bn = c$ to have solutions is when $c$ is the GCD (\href{https://artofproblemsolving.com/wiki/index.php?title=Greatest\_common\_divisor}{greatest common divisor}) of $a$ and $b$. Therefore, the answer is $gcd(300,210)=\boxed{\textbf{(C) }30}.$




\textbf{Solution 2 (Divisibility Analysis)}

Alternatively, note that $300p + 210g = 30(10p + 7g)$ is divisible by $30$ no matter what $p$ and $g$ are, so our answer must be divisible by $30$. Since we want the smallest integer, we can suppose that the answer is $30$ and go on from there. Note that three goats minus two pigs gives us $630 - 600 = 30$ exactly.  Since our supposition can be achieved, the answer is $\boxed{\textbf{(C) }30}$.

 



\textbf{Solution 3 (Simplifying the Problem)}

Let us simplify this problem.  Dividing by $30$, we get a pig to be:  $\frac{300}{30} =  10$, and a goat to be $\frac{210}{30}= 7$. It becomes evident that if you exchange $5$ pigs for $7$ goats, we get the smallest positive difference - $5\cdot 10 - 7\cdot 7 = 50-49 = 1$, since we can't make a non-integer with a linear combination of integers.  Since we originally divided by $30$, we need to multiply again, thus getting the answer  $1\cdot 30 = \boxed{\textbf{(C) }30}$.

 


	\item (\textit{AMC 10 10A 2006}) Circles with centers $A$ and $B$ have radii $3$ and $8$, respectively. A common internal tangent intersects the circles at $C$ and $D$, respectively. Lines $AB$ and $CD$ intersect at $E$, and $AE=5$. What is $CD$?


 
                    \begin{figure}[H]
                        \centering
                        \includegraphics[width=0.5\textwidth]{Problem/AMC_10_10A_2006_23_Problem_Figure_1.png}
                    \end{figure}
                    $\mathrm{(A)}\ 13\qquad\mathrm{(B)}\ \frac{44}{3}\qquad\mathrm{(C)}\ \sqrt{221}\qquad\mathrm{(D)}\ \sqrt{255}\qquad\mathrm{(E)}\ \frac{55}{3}$

 \textbf{Answer: }B



\textbf{Solution}


                    \begin{figure}[H]
                        \centering
                        \includegraphics[width=0.5\textwidth]{Solution/AMC_10_10A_2006_23_Solution_Figure_1.png}
                    \end{figure}
                    
$\triangle ADE \sim \triangle BCE$ because
 \[\angle AED \cong \angle BEC \text{ (vertical angles)}\]
 \[\angle ADE \cong \angle BCE \text{ (tangent perpendicularity)}\] 

 By the \href{https://artofproblemsolving.com/wiki/index.php?title=Pythagorean\_Theorem}{Pythagorean Theorem}, line segment $DE = 4$.  The sides are proportional, so 
 \[\frac{DE}{AD} = \frac{CE}{BC} \Rightarrow \frac{4}{3} = \frac{CE}{8}\Rightarrow CE = \frac{32}{3}\]

 And hence:
 \[CD = CE + DE = 4 + \frac{32}{3} = \boxed{\textbf{(B) }\frac{44}{3}}\]



	\item (\textit{AMC 10 10A 2006}) Centers of adjacent faces of a unit cube are joined to form a regular octahedron. What is the volume of this octahedron? 


 $\mathrm{(A)}\ \frac{1}{8}\qquad\mathrm{(B)}\ \frac{1}{6}\qquad\mathrm{(C)}\ \frac{1}{4}\qquad\mathrm{(D)}\ \frac{1}{3}\qquad\mathrm{(E)}\ \frac{1}{2}$

 \textbf{Answer: }B



\textbf{Solution}

We can break the octahedron into two \href{https://artofproblemsolving.com/wiki/index.php?title=Square\_pyramid}{square pyramids} by cutting it along a \href{https://artofproblemsolving.com/wiki/index.php?title=Plane}{plane} \href{https://artofproblemsolving.com/wiki/index.php?title=Perpendicular}{perpendicular} to one of its internal diagonals.
                    \begin{figure}[H]
                        \centering
                        \includegraphics[width=0.5\textwidth]{Solution/AMC_10_10A_2006_24_Solution_Figure_1.png}
                    \end{figure}
                    The cube has edges of length 1 so all edges of the regular octahedron have length $\frac{\sqrt{2}}{2}$.  Then the square base of the \href{https://artofproblemsolving.com/wiki/index.php?title=Pyramid}{pyramid} has area $\left(\frac{1}{2}\sqrt{2}\right)^2 = \frac{1}{2}$.We also know that the height of the pyramid is half the height of the cube, so it is $\frac{1}{2}$.  The volume of a pyramid with base area $B$ and height $h$ is $A=\frac{1}{3}Bh$ so each of the pyramids has volume $\frac{1}{3}\left(\frac{1}{2}\right)\left(\frac{1}{2}\right) = \frac{1}{12}$. The whole octahedron is twice this volume, so $\frac{1}{12} \cdot 2 = \boxed{\textbf{(B) }\frac{1}{6}}$.

 



\textbf{Solution 2}

The base of the Octahedron is exactly $\frac{1}{2}$ the base of the unit cube and the volume of the Octahedron is $\frac{1}{3}$ the size of the rectangular prism it occupies.

 Thus$\frac{1}{2} \cdot \frac{1}{3} = \boxed{\textbf{(B) }\frac{1}{6}}$.

 


	\item (\textit{AMC 10 10A 2006}) A bug starts at one vertex of a cube and moves along the edges of the cube according to the following rule. At each vertex the bug will choose to travel along one of the three edges emanating from that vertex. Each edge has equal probability of being chosen, and all choices are independent. What is the probability that after seven moves the bug will have visited every vertex exactly once? 


 $\mathrm{(A)}\ \frac{1}{2187}\qquad\mathrm{(B)}\ \frac{1}{729}\qquad\mathrm{(C)}\ \frac{2}{243}\qquad\mathrm{(D)}\ \frac{1}{81}\qquad\mathrm{(E)}\ \frac{5}{243}$

 \textbf{Answer: }C



\textbf{Solution 1}

Call this cube $ABCDEFGH$, with $A$ being the starting point.

 Because in $7$ moves the bug has to visit the other vertices, the bug cannot revisit any vertex.

 Therefore, starting at $A$, the bug has a $\frac{3}{3}$ chance of finding a good path to the next vertex, and call it $B$.

 Then, the bug has a $\frac{2}{3}$ chance of reaching a new vertex next. Call this $C$. $A, B,$ and $C$ are always on the same plane.

 Now, we split cases.

 Case $1$: The bug goes to the vertex $E$ opposite $A$ on the space diagonal with probability $\frac{1}{3}$.

 Then, the bug has to visit $D$ on the plane of $ABC$ last, as there is no way in and out from $D$.

 Therefore, there is only $1$ way out of $81$ to get to $D$ last.

 Therefore, there is a $\frac{1}{243} \cdot \frac{6}{9} = \frac{6}{2187}$ chance of finding a good path in this case.

 Case $2$: The bug goes to vertex $D$ on plane $ABC$ with a chance of $\frac{1}{3}$.

 The bug then has only $1$ way to go to a point $E$ on the opposite face, therefore having a $\frac{1}{3}$ probability.

 Then, the bug has a choice of two vertices on the face opposite to $ABCD$.

 This results in a $\frac{2}{3}$ probability of finding a good path to a point $F$.

 Then, there is only $1$ way out of $9$ to visit both other vertices on that face in $2$ moves.

 Multiply the probabilities for this case to get $\frac{2}{243} \cdot \frac{6}{9} = \frac{12}{2187}$.

 Add the probabilities of these two cases together to get $\frac{18}{2187} = \boxed{\textbf{(C) }\frac{2}{243}}.$




\textbf{Solution 2}


                    \begin{figure}[H]
                        \centering
                        \includegraphics[width=0.5\textwidth]{Solution/AMC_10_10A_2006_25_Solution_Figure_1.png}
                    \end{figure}
                    

 Let us count the good paths. The bug starts at an arbitrary \href{https://artofproblemsolving.com/wiki/index.php?title=Vertex}{vertex}, moves to a neighboring vertex ($3$ ways), and then to a new neighbor ($2$ more ways).  So, without loss of generality, let the cube have vertices $ABCDEFGH$ such that $ABCD$ and $EFGH$ are two opposite faces with $A$ above $E$ and $B$ above $F$.  The bug starts at $A$ and moves first to $B$, then to $C$.

 From this point, there are two cases. 

 Case $1$: the bug moves to $D$.  From $D$, there is only one good move available, to $H$.  From $H$, there are two ways to finish the trip, either by going $H \to G \to F \to E$ or $H \to E \to F \to G$.  So there are $2$ good paths in this case.

 Case $2$: the bug moves to $G$.

 Case $2a$: the bug moves $G \to H$.  In this case, there are $0$ good paths because it will not be possible to visit both $D$ and $E$ without double-visiting some vertex.

 Case $2b$: the bug moves $G \to F$.  There is $1$ good path in this case, $F \to E \to H \to D$.

 Thus, there are a total of $3\cdot 3 \cdot 2 = 18$ good paths.  There were $3^7 = 2187$ possible paths the bug could have taken, so the probability of a random path is good is $\frac{18}{2187} = \boxed{\textbf{(C) }\frac2{243}}$.

 



\textbf{Solution 3 (answer choices)}

As in Solution 1, the bug can move from its arbitrary starting vertex to a neighboring vertex in $3$ ways. After this, the bug can move to a new neighbor in $2$ ways (it cannot return to the first vertex). The total number of paths (as stated above) is $3^7$ or $2187$. Therefore, the probability of the bug following a good path is equal to $\frac{6x}{2187}$ for some positive integer $x$. The only answer choice which can be expressed in this form is $\boxed{\textbf{(C) }\frac2{243}}$.

 



\textbf{Solution 4}

Pick 1 vertex to start from. Notice that on all the odd moves the bug will move to 1 of 4 other vertices. Notice that on all even moves the bug must go to 1 of 3 other vertices (4 including the starting vertex). Thus, the bug will end on one of the odd vertices. There are 3 choices for the bug's first move from the starting vertex, and 2 choices for the bug's second move (it can't go back to the previous vertices). Then, notice that the bug has 3 choices of an odd vertex to end on not including the vertex the bug went to on move 1. Now pick 1 of the 3 ending vertices to end on. Notice that there are only 6 valid paths from start to finish (1 valid path after moves 1 and 2) to that point because all the moves are forced after the first 2 moves  (try it out). Thus, there are $3 \cdot 2 \cdot 3 = 18$ valid paths the bug can take. There are $3^7 = 2187$total paths, so the answer is 18/2187 = $\boxed{\textbf{(C) }\frac{2}{243}}$.

 



\textbf{Solution 5 (classical probability)}

Let the cube have vertices $ABCDEFGH$ such that $ABCD$ and $EFGH$ are two opposite faces with $A$ above $E$ and $B$ above $F$.  The bug starts at $A$ and every step has 3 choices so that there're $3^7 = 2187$ possible paths. We can split the case into odd and even steps. After an odd number of steps, the possible points that can be reached are $B$,$D$,$G$,$E$. Then, for even-numbered steps, the remaining points can be reached. Since he has taken seven steps, the end point must be one of the points in $B$,$D$,$G$,$E$. At this point, we classify and discuss the movement trajectory of the bug. 

 Case 1: It moves to another plane after completing one plane. If it starts on the upper plane, there are two possibilities for the first step, $B$ or $D$. The second and third steps are restricted to the upper plane, and the fourth step will move to the nearby point on another plane. These three steps have only one possibility each, so there is no need for classification. Then, when it reaches the lower plane, the fifth step has two possibilities, thus limiting the sixth and seventh steps as well. Proceeding in this way, the bug can also move from left to right or from front to back. Therefore, there are $2 \cdot 2 \cdot 3 = 12$ situations. 

 Case 2: Opposite to case 1, it can directly pass through to another point G on the other diagonal of the cube within the first three steps. There are three choices for the first step, two choices for the second step, and only one choice for the third step. Therefore, there are $3 \cdot 2 = 6$ situations for the first three steps. The subsequent steps will be restricted to a single path (try it yourself). 

 In summary, there are 12 + 6 = 18 situations in total. Therefore, the final probability is 18/2187 = $\boxed{\textbf{(C) }\frac2{243}}$.

 



\textbf{Solution 6}

We can tell that the total number of paths will be $3^6$ as the first and last vertex of the cube will be forced. We can now let the bottom left vertex of the front face of the cube be $A$ and, without loss of generality, assume the bug starts here. Now we can list out the possibilities for the bug to touch each vertex once after $7$ moves starting at $A$; we find $6$ favorable outcomes of the $729$, so we find the answer to be $\frac{6}{729} = \boxed{\textbf{(C) }\frac{2}{243}}$.

 


\end{enumerate}